% ══════════════════════════════════════════════════════════════════
%  CORRIGÉS — regroupés en fin d'ouvrage.
%  On corrige ici les exercices les plus exigeants (★★★) et les sujets type Bac.
% ══════════════════════════════════════════════════════════════════
\chapter{Corrigés des exercices}
\textit{Sont corrigés en détail ci-après les exercices les plus exigeants (★★★) et
l'ensemble des sujets type Baccalauréat de chaque chapitre. Les autres exercices,
plus directs, sont laissés au travail personnel de l'élève — l'essentiel étant de
chercher longuement avant de consulter une solution.}

% ──────────────────────────────────────────────────────────────────
\section{Chapitre 1 — Généralités sur la chimie organique}

\corrige{de l'exercice 6 (★★★)}
\begin{enumerate}[label=\arabic*)]
\item Masse de carbone : $m_C = m_{\ce{CO2}}\times\dfrac{12}{44} = 0{,}44\times\dfrac{12}{44} = \SI{0.12}{\gram}$.
Masse d'hydrogène : $m_H = m_{\ce{H2O}}\times\dfrac{2}{18} = 0{,}18\times\dfrac{2}{18} = \SI{0.02}{\gram}$.
Masse d'oxygène : $m_O = m - m_C - m_H = 0{,}30 - 0{,}12 - 0{,}02 = \SI{0.16}{\gram}$.
\item Quantités : $n_C = \dfrac{0{,}12}{12} = \SI{0.01}{\mol}$ ; $n_H = \dfrac{0{,}02}{1} = \SI{0.02}{\mol}$ ;
$n_O = \dfrac{0{,}16}{16} = \SI{0.01}{\mol}$.
On a $n_C : n_H : n_O = 1 : 2 : 1$, d'où le motif $\ce{CH2O}$ de masse molaire \SI{30}{\gram\per\mol}.
Comme $M = \SI{60}{\gram\per\mol} = 2\times30$, la formule brute est $\ce{C2H4O2}$.
\item Le composé est un acide ($\ce{C2H4O2}$ avec un groupe \ce{-COOH}) : c'est
l'\textbf{acide éthanoïque} (acide acétique) $\ce{CH3-COOH}$.
\end{enumerate}

\corrige{du sujet 1}
\begin{enumerate}[label=\arabic*)]
\item Quantité de $A$ vaporisé : $n = \dfrac{V}{V_m} = \dfrac{0{,}24}{24} = \SI{0.01}{\mol}$.
Masse molaire : $M = \dfrac{m}{n} = \dfrac{0{,}58}{0{,}01} = \SI{58}{\gram\per\mol}$.
\item $m_C = 1{,}32\times\dfrac{12}{44} = \SI{0.36}{\gram}$, soit $n_C = \SI{0.03}{\mol}$ ;
$m_H = 0{,}54\times\dfrac{2}{18} = \SI{0.06}{\gram}$, soit $n_H = \SI{0.06}{\mol}$ ;
$m_O = 0{,}58 - 0{,}36 - 0{,}06 = \SI{0.16}{\gram}$, soit $n_O = \SI{0.01}{\mol}$.
Ainsi $n_C : n_H : n_O = 3 : 6 : 1$ : la formule brute est $\ce{C3H6O}$ (vérification :
$M = 3\times12 + 6 + 16 = \SI{58}{\gram\per\mol}$).
\item Isomères de $\ce{C3H6O}$ : le propanal $\ce{CH3-CH2-CHO}$ (aldéhyde), la
propanone $\ce{CH3-CO-CH3}$ (cétone) et le prop-2-én-1-ol $\ce{CH2=CH-CH2-OH}$
(alcool insaturé).
\item $A$ réagit avec la 2,4-DNPH : il possède un groupe carbonyle $\ce{C=O}$ (aldéhyde
ou cétone). Mais il ne réduit pas la liqueur de Fehling : ce n'est pas un aldéhyde.
$A$ est donc une cétone : la \textbf{propanone} $\ce{CH3-CO-CH3}$.
\end{enumerate}

\corrige{du sujet 2}
\begin{enumerate}[label=\arabic*)]
\item Masse molaire : $M = 29\,d = 29\times1{,}93 \approx \SI{56}{\gram\per\mol}$.
\item À pression et température fixées, les volumes sont proportionnels aux quantités
de matière. Le volume de $\ce{CO2}$ formé valant quatre fois celui de $X$, une mole de
$X$ libère quatre moles de $\ce{CO2}$ : la molécule contient donc \textbf{4 atomes de carbone}.
\item Avec 4 carbones : $4\times12 = 48$, donc la masse d'hydrogène par mole est
$56 - 48 = \SI{8}{\gram\per\mol}$, soit 8 atomes d'hydrogène. La formule brute est
$\ce{C4H8}$. Cette formule de type $\ce{C_nH_{2n}}$ correspond à un composé
\textbf{insaturé} : un alcène (but-1-ène, but-2-ène $Z$ ou $E$, méthylpropène) ou un
cycloalcane (cyclobutane).
\end{enumerate}
